\documentclass[11pt]{article}
\usepackage[margin=2cm]{geometry}
\usepackage{hyperref}
\hypersetup{
    colorlinks=true,
    linkcolor=blue,
    citecolor=red,
    filecolor=magenta,      
    urlcolor=blue
    }
\usepackage{multirow}
\usepackage{amssymb,amsmath}
\usepackage{mathrsfs}
\usepackage[nameinlink,capitalize]{cleveref}

\usepackage{algorithm,algpseudocodex}
\renewcommand{\vec}[1]{\boldsymbol{#1}}
\newcommand{\intomega}[1]{\langle{#1}\rangle_\Omega}
\newcommand{\adjoint}[1]{{#1}^{\text{ad}}}
\newcommand{\xir}{\mathring{\widetilde{\xi}}}
\title{Lippmann Schwinger solver for linear elasticity}
\author{Eike Mueller, University of Bath}
\begin{document}
\maketitle
%%%%%%%%%%%%%%%%%%%%%%%%%%%%%%%%%%%%%%%%%%%%%%%%%%%%%%%%%%%%%%%%%%%%%%%%
\section{Problem description}
%%%%%%%%%%%%%%%%%%%%%%%%%%%%%%%%%%%%%%%%%%%%%%%%%%%%%%%%%%%%%%%%%%%%%%%%
We aim to solve the equations of linear elasticity
\begin{equation}
    \begin{aligned}
        \partial_j \sigma_{ij} &= 0 \\
        \sigma_{ij} &= C_{ijk\ell} \varepsilon_{k\ell}\qquad\text{with $\varepsilon_{k\ell} = \varepsilon^*_{k\ell} + \overline{\varepsilon}_{k\ell}$}\\
        \varepsilon^*_{k\ell} &= \frac{1}{2}\left(\partial_k u_\ell + \partial_\ell u_k\right)
    \end{aligned}\label{eqn:continuum_equations}
\end{equation}
in the domain $\Omega=[0,L_0]\times[0,L_0]\times[0,L_1]$
Here $C=C(x)$ is the given spatially varying elasticity tensor which has the following form for an isotropic material:
\begin{equation}
    C(x) = \lambda(x) \delta_{ij}\delta_{k\ell} + \mu(x) (\delta_{ik}\delta_{j\ell} + \delta_{i\ell}\delta_{jk})
\end{equation}
$\overline{\varepsilon}$ is a constant symmetric tensor representing the average strain and $u=u(x)$ is periodic such that $\int_\Omega u_i(x)\;dx = 0$ and $\int_\Omega \partial_i u_j(x)\;dx = 0$, i.e. $\int_\Omega \varepsilon^*_{ij}(x)\;dx=0$. This implies that $\frac{1}{|\Omega|}\int_\Omega \varepsilon_{ij}(x)\;dx=\overline{\varepsilon}_{ij}$
%%%%%%%%%%%%%%%%%%%%%%%%%%%%%%%%%%%%%%%%%%%%%%%%%%%%%%%%%%%%%%%%%%%%%%%%
\section{Discretisation}
%%%%%%%%%%%%%%%%%%%%%%%%%%%%%%%%%%%%%%%%%%%%%%%%%%%%%%%%%%%%%%%%%%%%%%%%
Consider the staggered grid in Fig. 1 of \cite{gelebart2020modified}, with the placement of variables shown at the voxel centres $\mathcal{C}$ and voxel corners (vertices) $\mathcal{V}$ in \cref{tab:placement}.
\begin{table}
    \caption{Placement of variables on grid}\label{tab:placement}
    \begin{center}
\begin{tabular}{ccc}
    \hline
location & variable & \\
\hline\hline
\multirow{2}{*}{voxel centres $\mathcal{C}$} & strain & $\varepsilon^*_{ij}=\frac{1}{2}(\partial_i u_j + \partial_j u_i)$ \\
& stress & $\sigma_{ij}=C_{ijk\ell} \varepsilon_{k\ell}$ \\\hline
\multirow{2}{*}{voxel corners $\mathcal{V}$} & displacement & $u_i$ \\
& stress divergence & $\partial_i \sigma_{ij}$ \\
\hline
\end{tabular}
\end{center}
\end{table}
Define the forward gradient $\vec{D}^+:\mathcal{V}\rightarrow \mathcal{C}$ of a variable $f\in \mathcal{V}$ placed at the voxel corners as
\begin{equation}
    \begin{aligned}
(D_0^+ f)_{abc} &= \frac{1}{4h_0}\Big[\left(f_{a+1,b,c}+f_{a+1,b+1,c}+f_{a+1,b,c+1}+f_{a+1,b+1,c+1}\right)\\
& \qquad -\;\;\left(f_{a,b,c}+f_{a,b+1,c}+f_{a,b,c+1}+f_{a,b+1,c+1} \right)\Big]
    \end{aligned}
\end{equation}
with corresponding expressions for $D_1^+ f$ and $D_2^+ f$.
The backward gradient $\vec{D}^-:\mathcal{C}\rightarrow \mathcal{V}$ of a variable $g\in \mathcal{C}$ placed at the voxel centres is given by
\begin{equation}
    \begin{aligned}
(D_0^- g)_{abc} &= \frac{1}{4h_0}\Big[\left(g_{a,b,c}+g_{a,b-1,c}+g_{a,b,c-1}+g_{a,b-1,c-1}\right)\\
& \qquad -\;\;\left(g_{a-1,b,c}+g_{a-1,b-1,c}+g_{a-1,b,c-1}+g_{a-1,b-1,c-1} \right)\Big]
    \end{aligned}
\end{equation}
Observe that $-\vec{D}^+\cdot \vec{D}^- = -\vec{D}^-\cdot \vec{D}^+$ is an approximation of the Laplace operator $\Delta=\partial_j\partial_j$.
%%%%%%%%%%%%%%%%%%%%%%%%%%%%%%%%%%%%%%%%%%%%%%%%%%%%%%%%%%%%%%%%%%%%%%%%
\section{Fourier expansion}
%%%%%%%%%%%%%%%%%%%%%%%%%%%%%%%%%%%%%%%%%%%%%%%%%%%%%%%%%%%%%%%%%%%%%%%%
The Fourier expansion of a general function $f$ is given by
\begin{equation}
    \begin{aligned}
\widehat{f}_{k_0,k_1,k_2} &= \sum_{n_0=0}^{N_0-1}\sum_{n_1=0}^{N_1-1}\sum_{n_2=0}^{N_2-1} f_{n_0,n_1,n_2}\exp\left[-\frac{2\pi i k_0 n_0}{N_0}-\frac{2\pi i k_1 n_1}{N_1}-\frac{2\pi i k_2 n_2}{N_2}\right] \\
f_{n_0,n_1,n_2} &= \frac{1}{N_0N_1N_2}\sum_{k_0=0}^{N_0-1}\sum_{k_1=0}^{N_1-1}\sum_{k_2=0}^{N_2-1} \widehat{f}_{k_0,k_1,k_2}\exp\left[\frac{2\pi i k_0 n_0}{N_0}+\frac{2\pi i k_1 n_1}{N_1}+\frac{2\pi  i k_2 n_2}{N_2}\right] 
    \end{aligned}
\end{equation}
For given $\vec{k}\in[0,N_0-1]\times[0,N_1-1]\times[0,N_2-1]$ define $\vec{\xi}\in[0,2\pi]^3$ with $\xi_i = \frac{2\pi k_i}{N_i}$, then the Fourier expansion can also be written as 
    \begin{xalignat}{2}
\widehat{f}_{\vec{\xi}} &= \sum_{n_0=0}^{N_0-1}\sum_{n_1=0}^{N_1-1}\sum_{n_2=0}^{N_2-1} f_{\vec{n}}\exp\left[-i\vec{\xi}\cdot \vec{n}\right], &
f_{\vec{n}} = \frac{1}{N_0N_1N_2}\sum_{k_0=0}^{N_0-1}\sum_{k_1=0}^{N_1-1}\sum_{k_2=0}^{N_2-1} \widehat{f}_{\vec{\xi}}\exp\left[i\vec{\xi}\cdot\vec{n}\right] 
    \end{xalignat}
It is then easy to see that in Fourier-space we have
\begin{equation}
\widehat{\vec{D}}^\pm = i \widetilde{\vec{\xi}} e^{\pm i \eta}
\end{equation}
with
\begin{equation}
\begin{aligned}
\widetilde{\xi}_0 &= \frac{2}{h_0}\sin\left(\frac{\xi_0}{2}\right)\cos\left(\frac{\xi_1}{2}\right)\cos\left(\frac{\xi_2}{2}\right)\\[1ex]
\widetilde{\xi}_1 &= \frac{2}{h_1}\cos\left(\frac{\xi_0}{2}\right)\sin\left(\frac{\xi_1}{2}\right)\cos\left(\frac{\xi_2}{2}\right)\\[1ex]
\widetilde{\xi}_2 &= \frac{2}{h_2}\cos\left(\frac{\xi_0}{2}\right)\cos\left(\frac{\xi_1}{2}\right)\sin\left(\frac{\xi_2}{2}\right)
\end{aligned}
\end{equation}
and
\begin{equation}
\eta = \frac{1}{2}\left(\xi_0+\xi_1+\xi_2\right).
\end{equation}
Observe that $\widehat{D}_i^+ \widehat{D}_j^- = \widehat{D}_i^- \widehat{D}_j^-= -\widetilde{\xi}_i \widetilde{\xi}_j$.
Consider the discretised homogeneous equations
\begin{xalignat}{2}
\sigma_{ij} &= C^0_{ijk\ell} D^+_k u_\ell + \tau_{ij}, & 
D^-_j \sigma_{ij} &= 0\label{eqn:discretised_equations}
\end{xalignat}
for a reference material with constant $C^0$. In Fourier-space this becomes
\begin{xalignat}{2}
\widehat{\sigma}_{ij} &= ie^{i\eta} C^0_{ijk\ell} \widetilde{\xi}_k \widehat{u}_\ell + \widehat{\tau}_{ij}, & 
ie^{-i\eta}\widetilde{\xi}_j \widehat{\sigma}_{ij} &= 0
\end{xalignat}
and therefore
\begin{equation}
K^0_{ik} \widehat{u}_k = ie^{-i\eta} \widehat{\tau}_{ij}\widetilde{\xi}_j
\end{equation}
with the acoustic tensor
\begin{equation}
    K^0_{ij} = C^0_{ikj\ell} \widetilde{\xi}_k\widetilde{\xi}_{\ell}.
\end{equation}
We also have
\begin{equation}
\widehat{\varepsilon}^*_{k\ell} = \frac{i}{2} e^{i\eta} \left(\widetilde{\xi}_k \widehat{u}_\ell+\widetilde{\xi}_\ell \widehat{u}_k\right)
\end{equation}
Putting everything together and following the derivation in the appendix of \cite{moulinec1998numerical} we find that for a homogeneous and isotropic reference material with
\begin{equation}
C^0_{ijk\ell} = \lambda^0 \delta_{ij}\delta_{k\ell} + \mu^0 \left(\delta_{ik}\delta_{j\ell}+\delta_{i\ell}\delta_{jk}\right)
\end{equation}
we have
\begin{equation}
    \begin{aligned}
\widehat{\varepsilon}^*_{k\ell ij} &= -\Gamma^0_{k\ell ij} \widehat{\tau}_{ij}\\
K^0_{ij} &= (\lambda^0+\mu^0)\widetilde{\xi}_i\widetilde{\xi}_j + \mu^0 |\widetilde{\vec{\xi}}|^2 \delta_{ij}\\
N^0_{ij} =\left((K^0)^{-1}\right)_{ij} &= \frac{1}{\mu^0 |\widetilde{\vec{\xi}}|^2}\left(\delta_{ij}-\frac{\widetilde{\xi}_i\widetilde{\xi}_j}{|\widetilde{\vec{\xi}}|^2}\frac{\lambda^0+\mu^0}{\lambda^0+2\mu^0}\right)\\
\widehat{\Gamma}^0_{k\ell ij} &= \frac{1}{4}\left(N^0_{\ell i}\widetilde{\xi}_j\widetilde{\xi}_k+N^0_{k i}\widetilde{\xi}_j\widetilde{\xi}_\ell + N^0_{\ell j}\widetilde{\xi}_i\widetilde{\xi}_k+N^0_{kj}\widetilde{\xi}_i\widetilde{\xi}_\ell\right)\\
&= \frac{1}{\mu^0}\left(\frac{1}{4}\left(\delta_{\ell i} \mathring{\widetilde{\xi}}_j\mathring{\widetilde{\xi}}_k+\delta_{ki} \mathring{\widetilde{\xi}}_j\mathring{\widetilde{\xi}}_\ell+\delta_{\ell j} \mathring{\widetilde{\xi}}_i\mathring{\widetilde{\xi}}_k+\delta_{k j} \mathring{\widetilde{\xi}}_i\mathring{\widetilde{\xi}}_\ell\right)-\frac{\lambda^0+\mu^0}{\lambda^0+2\mu^0}\mathring{\widetilde{\xi}}_i\mathring{\widetilde{\xi}}_j\mathring{\widetilde{\xi}}_k\mathring{\widetilde{\xi}}_\ell\right)\\
&= \frac{1}{\mu^0}\left( \widehat{A}_{k\ell ij} - \frac{\lambda^0+\mu^0}{\lambda^0+2\mu^0} \widehat{B}_{k\ell ij}\right)
    \end{aligned}\label{eqn:fourier_expression}
\end{equation}
with
\begin{equation}
\mathring{\widetilde{\xi}}_i = \begin{cases}
0 & \text{if $\|\widetilde{\vec{\xi}}\|=0$}\\
\widetilde{\xi}_i/\|\widetilde{\vec{\xi}}\| & \text{otherwise}
\end{cases}
\end{equation}
where $\|\cdot\|$ denotes the Euclidean two-norm and we defined
\begin{equation}
\begin{aligned}
\widehat{A}_{k\ell ij} &= \frac{1}{4}\left(\delta_{\ell i} \mathring{\widetilde{\xi}}_j\mathring{\widetilde{\xi}}_k+\delta_{ki} \mathring{\widetilde{\xi}}_j\mathring{\widetilde{\xi}}_\ell+\delta_{\ell j} \mathring{\widetilde{\xi}}_i\mathring{\widetilde{\xi}}_k+\delta_{k j} \mathring{\widetilde{\xi}}_i\mathring{\widetilde{\xi}}_\ell\right)\\
\widehat{B}_{k\ell ij} &= \mathring{\widetilde{\xi}}_i\mathring{\widetilde{\xi}}_j\mathring{\widetilde{\xi}}_k\mathring{\widetilde{\xi}}_\ell
\end{aligned}
\end{equation}
If we use Voigt's notation \cite{voigt1928lehrbuch} and the convention used in the AMITEX code to arrange the six independent components of the symmetric $3\times 3$ tensors $\widehat{\varepsilon}$ and $\widehat{\tau}$ into vectors
\begin{equation}
    \begin{aligned}
    \widehat{\vec{\varepsilon}}^* &=
    \begin{pmatrix}
        \widehat{\varepsilon}^*_{0} &
        \widehat{\varepsilon}^*_{1} &
        \widehat{\varepsilon}^*_{2} &
        \widehat{\varepsilon}^*_{3} &
        \widehat{\varepsilon}^*_{4} &
        \widehat{\varepsilon}^*_{5}
    \end{pmatrix}:= 
    \begin{pmatrix}
        \widehat{\varepsilon}^*_{00} &
        \widehat{\varepsilon}^*_{11} &
        \widehat{\varepsilon}^*_{22} &
        \widehat{\varepsilon}^*_{01} &
        \widehat{\varepsilon}^*_{02} &
        \widehat{\varepsilon}^*_{12}
    \end{pmatrix}\\
    \widehat{\vec{\tau}} &= 
    \begin{pmatrix}
        \widehat{\tau}_{0} &
        \widehat{\tau}_{1} &
        \widehat{\tau}_{2} &
        \widehat{\tau}_{3} &
        \widehat{\tau}_{4} &
        \widehat{\tau}_{5}
    \end{pmatrix}= 
    \begin{pmatrix}
        \widehat{\tau}_{00} &
        \widehat{\tau}_{11} &
        \widehat{\tau}_{22} &
        \widehat{\tau}_{01} &
        \widehat{\tau}_{02} &
        \widehat{\tau}_{12}
    \end{pmatrix}
\end{aligned}
\end{equation}
we can write
\begin{equation}
\widehat{\vec{\varepsilon}}^* = -\frac{1}{\mu^0}\left(\widehat{\vec{\varepsilon}}^{*(A)} - \frac{\lambda^0+\mu^0}{\lambda^0+2\mu^0} \widehat{\vec{\varepsilon}}^{*(B)}\right)
\end{equation}
with 
\begin{equation}
\begin{aligned}    
    \widehat{\vec{\varepsilon}}^{*(A)} &= \begin{pmatrix}
    \xir_0^2 \widehat{\tau}_0 + \xir_0(\xir_2\widehat{\tau}_4+\xir_1\widehat{\tau}_3)\\[2ex]
    \xir_1^2 \widehat{\tau}_1 + \xir_1(\xir_2\widehat{\tau}_5+\xir_0\widehat{\tau}_3)\\[2ex]
    \xir_2^2 \widehat{\tau}_2 + \xir_2(\xir_1\widehat{\tau}_5+\xir_0\widehat{\tau}_4)\\[2ex]
    \frac{1}{2}\left(\xir_0\xir_1(\widehat{\tau}_0+\widehat{\tau}_1)+(\xir_0^2+\xir_1^2)\widehat{\tau}_3+\xir_2(\xir_0\widehat{\tau}_5+\xir_1\widehat{\tau}_4)\right)\\[2ex]
    \frac{1}{2}\left(\xir_0\xir_2(\widehat{\tau}_0+\widehat{\tau}_2)+(\xir_0^2+\xir_2^2)\widehat{\tau}_4+\xir_1(\xir_0\widehat{\tau}_5+\xir_2\widehat{\tau}_3)\right)\\[2ex]
    \frac{1}{2}\left(\xir_1\xir_2(\widehat{\tau}_1+\widehat{\tau}_2)+(\xir_1^2+\xir_2^2)\widehat{\tau}_5+\xir_0(\xir_1\widehat{\tau}_4+\xir_2\widehat{\tau}_3)\right)
    \end{pmatrix}
\end{aligned}
\end{equation}
and
\begin{equation}
\begin{aligned}    
    \widehat{\vec{\varepsilon}}^{*(B)} &= 
    \left(
    \xir_0^2\widehat{\tau}_0+ \xir_1^2\widehat{\tau}_1+ \xir_2^2\widehat{\tau}_2+ 2(\xir_0\xir_1\widehat{\tau}_3+ \xir_0 \xir_2\widehat{\tau}_4+\xir_1\xir_2\widehat{\tau}_5)
    \right)\begin{pmatrix}
    \xir_0^2 \\[2ex] \xir_1^2 \\[2ex] \xir_2^2 \\[2ex] \xir_0\xir_1 \\[2ex] \xir_0 \xir_2 \\[2ex] \xir_1\xir_2
    \end{pmatrix}^\top 
\end{aligned}
\end{equation}
where the expression for $\widehat{\vec{\varepsilon}}^{*(A)}=\widehat{A}\widehat{\vec{\tau}}$ follows from the $6\times 6$ matrix representation of $\widehat{A}_{k\ell ij}$ in Voigt's notation
\begin{equation}
\widehat{A} = \begin{pmatrix}
    \xir_0^2 & 0 & 0 & \xir_0\xir_1 & \xir_0\xir_2 & 0 \\[2ex]
    0 & \xir_1^2 & 0 & \xir_0\xir_1 & 0 & \xir_1\xir_2\\[2ex]
    0 & 0 & \xir_2^2 & 0 & \xir_0\xir_2 & \xir_1\xir_2\\[2ex]
    \frac{1}{2}\xir_0\xir_1 & \frac{1}{2}\xir_0\xir_1 & 0 & \frac{1}{2}(\xir_0^2+\xir_1^2) & \frac{1}{2}\xir_1\xir_2 & \frac{1}{2}\xir_0\xir_2\\[2ex]
    \frac{1}{2}\xir_0\xir_2 & 0 & \frac{1}{2}\xir_0\xir_2 & \frac{1}{2}\xir_1\xir_2 & \frac{1}{2}(\xir_0^2+\xir_2^2) & \frac{1}{2}\xir_0\xir_1\\[2ex]
    0 & \frac{1}{2}\xir_1\xir_2 & \frac{1}{2}\xir_1\xir_2 & \frac{1}{2}\xir_0\xir_2 & \frac{1}{2}\xir_0\xir_1 & \frac{1}{2}(\xir_1^2+\xir_2^2)
    \end{pmatrix}
\end{equation}
%%%%%%%%%%%%%%%%%%%%%%%%%%%%%%%%%%%%%%%%%%%%%%%%%%%%%%%%%%%%%%%%%%%%%%%%
\section{Lippmann Schwinger iteration}
%%%%%%%%%%%%%%%%%%%%%%%%%%%%%%%%%%%%%%%%%%%%%%%%%%%%%%%%%%%%%%%%%%%%%%%%
To solve \eqref{eqn:continuum_equations}, observe that the discretisation of these equations is given by \eqref{eqn:discretised_equations} with
\begin{equation}
    \begin{aligned}
    \tau_{ij} &= (C_{ijk\ell} - C^0_{ijk\ell}) \varepsilon_{k\ell} \\
    &= (\lambda-\lambda^0) \varepsilon_{kk} \delta_{ij} + 2(\mu-\mu^0)\varepsilon_{ij}.
    \end{aligned}
\end{equation}
We can write $\tau_{ij} = C_{ijk\ell}  \varepsilon_{k\ell}$ explicitly as
\begin{equation}
\begin{aligned}
\tau_{00} &= 2\mu\varepsilon_{00} + \lambda\left(\varepsilon_{00}+\varepsilon_{11}+\varepsilon_{22}\right) \\
\tau_{11} &= 2\mu\varepsilon_{11} + \lambda\left(\varepsilon_{00}+\varepsilon_{11}+\varepsilon_{22}\right)\\
\tau_{22} &= 2\mu\varepsilon_{22} + \lambda\left(\varepsilon_{00}+\varepsilon_{11}+\varepsilon_{22}\right)\\
\tau_{01} &= 2\mu\varepsilon_{01}\\
\tau_{02} &= 2\mu\varepsilon_{02}\\
\tau_{12} &= 2\mu\varepsilon_{12}.
\end{aligned}\label{eqn:tau_computation}
\end{equation}
\begin{algorithm}
    \caption{Lippmann Schwinger iteration: simplest formulation}
\begin{algorithmic}[1]
    \State{Set $\varepsilon=E$}
    \For{$n=0,1,\dots$}
    \State{Compute stress $\sigma_{ij}=C_{ijk\ell}\varepsilon_{k\ell}$ according to \eqref{eqn:tau_computation}}
    \State{Check convergence based on $D^-_j \sigma_{ij}$}
    \State{Compute stress correction $\tau_{ij}=\sigma_{ij}-C^0_{ijk\ell} \varepsilon_{k\ell}$}
    \State{Compute $\widehat{\tau}_{ij} = \mathscr{F}[\tau_{ij}]$}\Comment{Fourier-transform}
    \State{Set $\widehat{\varepsilon}^*_{k\ell} = -\widehat{\Gamma}^0_{k\ell ij} \widehat{\tau}_{ij}$}\Comment{Solve in Fourier space}
    \State{Update $\varepsilon_{k\ell} = \mathscr{F}^{-1}[\widehat{\varepsilon}^*_{k\ell}]+E\overline{\varepsilon}_{k\ell}$}\Comment{Inverse Fourier-transform}        
    \EndFor
\State{\Return $\varepsilon$}
\end{algorithmic}
\end{algorithm}
\begin{algorithm}
    \caption{Lippmann Schwinger iteration: residual formulation}
\begin{algorithmic}[1]
    \State{Set $\varepsilon=E$}
    \For{$n=0,1,\dots$}
    \State{Compute stress $\sigma_{ij}=C_{ijk\ell}\varepsilon_{k\ell}$ according to \eqref{eqn:tau_computation}}
    \State{Check convergence based on $D^-_j \sigma_{ij}$}    
    \State{Compute $\widehat{\sigma}_{ij} = \mathscr{F}[\sigma_{ij}]$}\Comment{Fourier-transform}
    \State{Compute $\widehat{r}_{k\ell} = -\widehat{\Gamma}^0_{k\ell ij} \widehat{\sigma}_{ij}$}\Comment{Solve in Fourier space}
    \State{Compute $r_{ij} = \mathscr{F}^{-1}[\widehat{r}_{ij}]$}\Comment{Inverse Fourier-transform}
    \State{Update $\varepsilon_{k\ell} \mapsto \varepsilon_{k\ell} + r_{k\ell}$}
    \EndFor
\State{\Return $\varepsilon$}
\end{algorithmic}
\end{algorithm}
\subsection{Anderson acceleration}
The Lippmann-Schwinger algorithm with Anderson acceleration \cite{wicht2021anderson} is shown in \cref{alg:lippmann_schwinger_anderson}. It requires additional storage of $2\times(d+1)$ tensors for the state $\varepsilon^{(s)}$ and residuals $r^{(s)}$, as well as a $(d+1)\times (d+1)$ matrix and vectors $u,v\in\mathbb{R}^{d+1}$.
\begin{algorithm}
    \caption{Lippmann Schwinger iteration: Anderson acceleration with depth $d$}\label{alg:lippmann_schwinger_anderson}
\begin{algorithmic}[1]
    \State{Set $\varepsilon=E$}
    \State{Initialise $\varepsilon^{(0)} = \varepsilon$, $\varepsilon^{(j)} = 0$ for $j=1,2,\dots,d$}
    \State{Initialise $r^{(0)} = 0$ for $s=0,1,2,\dots,d$}
    \State{Initialise $(d+1)\times (d+1)$ identity matrix $A=\mathbb{I}$}
    \For{$n=0,1,\dots$}
    \State{Compute stress $\sigma_{ij}=C_{ijk\ell}\varepsilon_{k\ell}$ according to \eqref{eqn:tau_computation}}
    \State{Check convergence based on $D^-_j \sigma_{ij}$}    
    \State{Compute $\widehat{\sigma}_{ij} = \mathscr{F}[\sigma_{ij}]$}\Comment{Fourier-transform}
    \State{Compute $\widehat{r}_{k\ell} = \widehat{\Gamma}^0_{k\ell ij} \widehat{\sigma}_{ij}$}\Comment{Solve in Fourier space}
    \State{Set $m_k=\min\{n,d\}$}
    \For{$s=m_k,m_k-1,\dots,2,1$}
    \State{Set $r^{(s)}\gets r^{(s-1)}$}
    \For{$t=m_k,m_k-1,\dots,2,1$}
    \State{Set $A_{st}\gets A_{s-1,t-1}$}
    \EndFor
    \EndFor
    \State{Compute $r^{(0)}_{ij} = \mathscr{F}^{-1}[\widehat{r}_{ij}]$}\Comment{Inverse Fourier-transform}
    \For{$s=0,1,2,\dots,m_k$}
    \State{Set $A_{0s} = A_{s0}\gets r^{(0)}_{ij}r^{(s)}_{ij}$}
    \EndFor
    \State{Solve $Av=u$ with $u=(1,1,\dots,1)\in\mathbb{R}^d$}
    \State{Set $\alpha = v/(u^\top v)$}
    \State{Set $\varepsilon_{ij} = \sum_{s=0}^{d} \alpha_s (\varepsilon^{(s)}_{ij}-r^{(s)}_{ij})$}
    \For{$s=m_k,m_k-1,\dots,2,1$}
    \State{Set $\varepsilon^{(s)}\gets \varepsilon^{(s-1)}$}
    \EndFor
    \State{Set $\varepsilon_{ij} \gets \varepsilon_{ij}^{(0)}$}
    \EndFor
\State{\Return $\varepsilon$}
\end{algorithmic}
\end{algorithm}
\bibliographystyle{unsrt}
\bibliography{linear_elasticity}
\end{document}